% GENERATED_BY: oc_core_monograph_translation_draft_pipeline_v1.py
% AI_ASSISTED_TRANSLATION_DRAFT: TRUE
% THEOREM_REVIEW_STATUS: PENDING
% THEOREM_REVIEW_MODE: FORMAL_AI_GATE__CODEX_CLI_CHATGPT
% THEOREM_REVIEW_REVIEWER_ID:
% THEOREM_REVIEW_AUTHORITY_CLASS:
% THEOREM_REVIEWED_AT:
% SOURCE_LANGUAGE: EN
% TARGET_LANGUAGE: RU
% SOURCE_EN_REF: content/m_spaces/mspaces_master.tex
% ================================================================
% ==== FILE: content/m_spaces/mspaces_master.tex
% ================================================================

% ==============================
%  Ontology of Continua — Core
%  M-Spaces: Master Overview
%  FULL MODULE — FINAL
% ==============================

\section{MSpaces}
\label{sec:mspaces-overview}

Семейство мета-пространств $M_0, M_1, \dots, M_{12}$ задаёт
\emph{структуру допустимости} для континуумов $K_0\dots K_{12}$.
Каждый $K_i$ может существовать только если его полная конфигурация
$\Omega(K_i)$ строго лежит внутри допустимой области соответствующего
мета-пространства $M_i$:
\[
\Omega(K_i) \subseteq \Omega(M_i).
\]

M-пространства выполняют три универсальные функции:

\begin{enumerate}
    \item \textbf{Ограничения допустимости:}
    они определяют, какие структурные конфигурации, пороги и потоки допустимы.
    Они служат глобальным регулятором континуумов.

    \item \textbf{Размерностное направление:}
    $M_{i+1}$ определяет, может ли $K_i$ формировать новые оси и
    переходить в континуум более высокой размерности.

    \item \textbf{Структурный надзор:}
    каждое мета-пространство обеспечивает глобальную совместимость операторов
    таких как $\Psi$, $\Phi$, $\Lambda$, $U$, и $\Chi$, и гарантирует,
    что континуумы удовлетворяют универсальным законам эволюции, коллапса и
    ветвления.
\end{enumerate}

M-пространства не эволюционируют «внутри» иерархии $K$; вместо
этого они образуют параллельную надзорную структуру.
Взаимосвязь можно визуализировать как:
\[
K_i \hookrightarrow M_i,\qquad
M_i \Rightarrow \text{ограничения на } K_i.
\]

Переход $K_i \to K_{i+1}$ возможен только если:
\[
A(K_{i+1}) \subseteq A(M_{i+1}), \quad
P(K_{i+1}) \subseteq P(M_{i+1}),\quad
\Theta(K_{i+1}) \leq \Theta(M_{i+1}).
\]

Каждое $M_i$ расширяет размерность допустимых структур по сравнению с
$M_{i-1}$, формируя надструктуру, в которой может возникнуть
следующий уровень континуума.
